% Generated draft: requires mathematical and source review.
\documentclass[11pt]{article}
\usepackage[margin=25mm]{geometry}
\usepackage{fontspec}
\setmainfont{Latin Modern Roman}
\usepackage{amsmath,amssymb,mathtools,unicode-math}
\setmathfont{Latin Modern Math}
\usepackage{enumitem,needspace,xurl,hyperref}
\setlength{\parindent}{0pt}
\setlength{\parskip}{5pt}
\setlength{\emergencystretch}{3em}
\newcommand{\N}{\mathbb{N}}
\newcommand{\Z}{\mathbb{Z}}
\newcommand{\Q}{\mathbb{Q}}
\newcommand{\R}{\mathbb{R}}
\newcommand{\C}{\mathbb{C}}
\newcommand{\F}{\mathbb{F}}
\newcommand{\A}{\mathbb{A}}
\newcommand{\PP}{\mathbb{P}}
\newcommand{\E}{\mathbb{E}}
\newcommand{\eps}{\varepsilon}
\newcommand{\dd}{\,\mathrm{d}}
\begin{document}
\section*{UnsolvedMath ID: 30002896}
\textbf{Source problem code:} catalogue entry
\subsection{Definitions and mathematical statement}
A Markov shift is a shift on a path space equipped with a Markov law. Its specification requires an index set, a measurable state space, and transition and initial or marginal laws. A measurable transformation \(T\) on a measure space \((X,\mu)\) is nonsingular if \(\mu(T^{-1}A)=0\) whenever \(\mu(A)=0\). The supplied background describes the intended Markov shifts as nonsingular, but gives no further hypotheses on them.

A Bernoulli shift is the shift on a product space whose coordinates are independent and identically distributed under the product measure. Its precise specification requires an indexing action and a base probability space. In the ordinary probabilistic sense, a coupling of probability measures \(\mu\) on \(X\) and \(\nu\) on \(Y\) is a probability measure \(\lambda\) on \(X\times Y\) satisfying
\[
(\pi_X)_*\lambda=\mu,\qquad (\pi_Y)_*\lambda=\nu,
\]
where \(\pi_X\) and \(\pi_Y\) are the coordinate projections. This definition alone does not specify the dynamical meaning of ``representable as a coupling to'' in the source.

The supplied question is: Is each Markov shift \(\rho(x_n)\) representable as a coupling to a Bernoulli shift? The background qualifies the intended shifts as nonsingular. The universal quantifier ranges over the intended class of such Markov shifts, and the Bernoulli shift may depend on the Markov shift. However, the source does not define that class, the notation \(\rho(x_n)\), or the required coupling and representation conditions. Thus the question is retained as an unresolved formulation; a fully specified mathematical assertion cannot be extracted without those definitions.
\subsection{Short English statement}
The question asks whether every intended nonsingular Markov shift \(\rho(x_n)\) can be represented through a coupling with some Bernoulli shift. The class of shifts and the meaning of the required representation are unspecified, so this statement requires clarification and human mathematical review.
\subsection{Sources}
\begin{itemize}
\item Supplied local UnsolvedMath registry record: permanent ID \(30002896\), title ``Coupling Markov Shifts to Bernoulli Shifts,'' Oberwolfach Reports Open Problems collection, exact source problem code \texttt{catalogue entry}. The record gives proposed year 2015 and names no proposer.
\item Supplied report citation: \emph{Free Probability Theory (2016)}, Oberwolfach Reports, DOI 10.4171/owr/2015/28, pp.\ 1571--1629. \url{https://doi.org/10.4171/owr/2015/28}. The background describes Oberwolfach Reports as documenting mathematical workshops and their open-problem discussions. The supplied literature-source label is ``Oberwolfach report 28.''
\item Supplied literature source: \emph{All properly ergodic Markov chains over a free group are orbit equivalent}. \url{https://arxiv.org/abs/1706.09483}.
\item Literature assessment dated exactly \texttt{2026-08-21}: the supplied curation labels the problem open and states that the report poses the coupling assertion for the specified nonsingular Markov shifts, while later classification and orbit-equivalence results do not settle the general assertion. This is supplied dated curation, not an independently verified assessment. The background expressly requests independent verification of time-sensitive status and cited claims.
\item The registry labels the statement \texttt{corrected\_verified}, explaining that only safe notation, typography, and TeX-delimiter cleanup was applied and that hypotheses, quantifiers, and subquestions were preserved. This supplied label does not establish mathematical correctness or resolve the missing definitions.
\end{itemize}
\end{document}
